\documentclass[10pt]{article}
\usepackage{lmodern}
\usepackage[letterpaper,margin=0.8in]{geometry}
\usepackage{amsmath,amssymb,booktabs,array,longtable}
\usepackage[hidelinks]{hyperref}
\usepackage{fancyvrb}
\hypersetup{pdftitle={Recovery-Set Contracts for Cryptographic Storage Migration: Supplementary Material},pdfauthor={Ishan Teredesai and Tushar Rahangdale}}
\title{Recovery-Set Contracts for Cryptographic Storage Migration\\\large Supplementary Material and Reproduction Guide}
\author{Ishan Teredesai \qquad Tushar Rahangdale}
\date{8 November 2026}
\begin{document}
\maketitle
\section{Scope and reviewer access}
This artifact contains a bounded reference model, a process-interruption harness,
and a two-worker concurrent-admission harness. The implementations use SQLite,
AES-256-GCM, and an unprotected local JSON key file. That file is a test fixture,
not a secure KMS. There is no ML-KEM/hybrid-KEM implementation, QRNG, industrial
case study, power-loss experiment, performance benchmark, or cryptographic proof.
All data are synthetic. Counts describe selected cases, not failure probabilities.

The private repository is \url{https://github.com/titcsl/storage-recovery-contracts}.
A private URL alone does not grant reviewer access. The complete ZIP includes
source, raw observations, traces, manifests, this guide, and current LaTeX/PDFs;
it can be transferred through the venue's permitted artifact mechanism without
repository access. No public access or reviewer invitation is claimed.

\section{Reproduction commands}
Use Python 3.10+ on Linux (Windows: WSL/Ubuntu). Unzip the complete artifact and
run commands from its \texttt{recovery-contracts/} directory. The model needs
only the standard library; both storage harnesses require \texttt{cryptography}.
\begin{Verbatim}[fontsize=\small]
python3 -m venv .venv
source .venv/bin/activate
python -m pip install -r requirements.txt
# One command for each experiment:
python3 experiment/model.py --out results
python3 experiment/crash_harness.py --out results
python3 experiment/concurrency_harness.py --out results
# Validate regression tests and consistency of recorded results:
python3 -m unittest discover -s experiment -p 'test_*.py' -v
python3 experiment/verify_results.py --results results
# Or run the entire sequence:
bash run_experiments.sh
\end{Verbatim}
Commands replace files in \texttt{results/}. Preserve the delivered directory
and rerun in a separate copy to compare environments. A nonzero exit is a failed
run; do not treat old files left in the directory as a new successful run.

\subsection*{Recorded successful environment}
\begin{center}\begin{tabular}{@{}ll@{}}
\toprule
Python & 3.12.14 (Clang 22.1.3; build dated 25 August 2026)\\
SQLite / cryptography & 3.53.1 / 50.0.2\\
Platform & Linux 6.18.57, x86-64, glibc 2.39\\
Concurrent run completed & 8 November 2026, 06:49:15 UTC\\
Scheduling and lock & Two spawned workers; IPC gates; local \texttt{fcntl.flock}\\
\bottomrule
\end{tabular}\end{center}
The JSON manifests preserve the full runtime strings, source SHA-256 values,
configuration bounds, and failure models. No cloud services are used by the code.

\newpage
\section{Finite model and all eighteen configurations}
State $S=(a,L,B,P,K)$ contains active generation, committed live references,
optional snapshot, volatile pending targets, and available keys. Initially live
and snapshot references are zero, the only key is zero, and no work is pending.
The invariant concerns live and retained-snapshot references. The retirement
guard additionally protects pending targets to preserve prepared work.

\begin{center}\begin{tabular}{@{}p{2.4cm}p{6.1cm}p{7cm}@{}}
\toprule
Transition & Precondition & Effect\\
\midrule
Activate & Next generation is within the bound. & Durably add key; change active generation.\\
Prepare $i$ & No pending target; source key exists; live generation differs from active. & Set volatile target to active generation.\\
Commit $i$ & Target key is available. & Publish target and clear pending entry.\\
Snapshot & No retained snapshot; live keys available. & Copy committed live references.\\
Release & Snapshot exists; authorized expiry assumed. & Remove that retention obligation, without migrating it.\\
Crash & Pending work exists. & Discard pending work, retain durable state.\\
Retire $k$ & Non-active key with no dependency under the selected rule. & Remove key from available inventory.\\
\bottomrule
\end{tabular}\end{center}
Breadth-first exploration ends only at a fixed point. Counts include edges to
previously visited states. Intent-only substitutes pending targets for live
references and ignores snapshots. Live-only checks committed live references.
Recovery-set checks live, snapshot, and pending references. These first two
policies are constructed negative controls, not reported product vulnerabilities.

\begin{center}\small
\begin{tabular}{@{}rrlrrrrr@{}}
\toprule
Records & Generations & Policy & States & Edges & Missing L & Missing B & Profile gap\\
\midrule
1 & 2 & intent-only & 16 & 31 & 4 & 3 & 2 \\
1 & 2 & live-only & 12 & 23 & 0 & 1 & 2 \\
1 & 2 & recovery-set & 11 & 21 & 0 & 0 & 1 \\
1 & 3 & intent-only & 76 & 202 & 26 & 20 & 14 \\
1 & 3 & live-only & 52 & 137 & 0 & 7 & 10 \\
1 & 3 & recovery-set & 43 & 111 & 0 & 0 & 5 \\
2 & 2 & intent-only & 48 & 115 & 16 & 13 & 6 \\
2 & 2 & live-only & 32 & 79 & 0 & 3 & 6 \\
2 & 2 & recovery-set & 29 & 73 & 0 & 0 & 3 \\
2 & 3 & intent-only & 500 & 1,642 & 230 & 202 & 126 \\
2 & 3 & live-only & 288 & 955 & 0 & 55 & 82 \\
2 & 3 & recovery-set & 217 & 713 & 0 & 0 & 39 \\
3 & 2 & intent-only & 154 & 448 & 52 & 45 & 14 \\
3 & 2 & live-only & 102 & 319 & 0 & 7 & 14 \\
3 & 2 & recovery-set & 95 & 305 & 0 & 0 & 7 \\
3 & 3 & intent-only & 3,682 & 14,499 & 1,814 & 1,766 & 926 \\
3 & 3 & live-only & 1,996 & 8,031 & 0 & 403 & 610 \\
3 & 3 & recovery-set & 1,495 & 6,009 & 0 & 0 & 293 \\
\bottomrule\end{tabular}\end{center}
Missing L/B counts states with unavailable live/backup keys; the columns can
overlap. Profile gap means live references are target-labeled but a retained
snapshot remains legacy. The guarded configurations total 1,890 states and
7,232 edges, counted separately, with no missing-key state. The largest has
293 profile-gap states. None of these quantities is a deployment frequency.

Release is an authorized retention action, not snapshot migration or secure
erasure. A guard may preserve reference availability while blocking retirement
indefinitely. Labels do not establish actual post-quantum protection.

\newpage
\section{Process interruption and retained snapshots}
Each record receives a fresh 32-byte DEK. Payload encryption and DEK wrapping
use AES-256-GCM with separate random 12-byte nonces. Payload associated data binds
record identity and content version; wrapping associated data binds identity
and epoch. The key file uses file synchronization and atomic replacement;
SQLite uses rollback journaling with \texttt{synchronous=FULL}. These are separate
durability domains. The test stops a process, not the kernel or power supply.

The controller provisions generation 1, prepares replacement envelopes, attempts
retirement, applies a SQL transaction, commits, and attempts retirement again.
A pre-rotation snapshot remains at generation 0. The parent restores both copies
using the available key file and checks plaintext bytes against generated input.
\begin{center}\begin{tabular}{@{}p{6.7cm}rrr@{}}
\toprule
Cut/completion point & Intent-only L/B & Live-only L/B & Recovery-set L/B\\
\midrule
After durable provisioning & 0/0 & 0/0 & 0/0\\
After envelope preparation & 0/0 & 0/0 & 0/0\\
After first retirement attempt & 6/6 & 0/0 & 0/0\\
After SQL updates, before commit & 6/6 & 0/0 & 0/0\\
After database commit & 0/6 & 0/0 & 0/0\\
After final retirement attempt & 0/6 & 0/6 & 0/0\\
Normal completion & 0/6 & 0/6 & 0/0\\
\bottomrule\end{tabular}\end{center}
Each cell reports failure counts across six datasets: one/eight records at
512 B, 16 KiB, or 256 KiB each. Every policy/cut/dataset runs once. There are
126 cases, with 108 child terminations at exit code 93 and 18 normal exits.
Columns may overlap; these are deterministic coverage outcomes.

\begin{center}\begin{tabular}{@{}lrrrr@{}}
\toprule
Policy & Cases & Live failures & Backup failures & Both restored\\
\midrule
Intent-only & 42 & 12 & 30 & 12\\
Live-only & 42 & 0 & 12 & 30\\
Recovery-set & 42 & 0 & 0 & 42\\
\bottomrule\end{tabular}\end{center}
An unavailable wrapping key counts as a failed restore. An authentication error,
wrong row count, unexpected child exit, or plaintext mismatch aborts the run.
The harness records observed restoration counts before checking its expected
control behavior; it does not substitute expected counts for decryption results.
The separate concurrent harness has no expected-result function.

\subsection*{Minimal model witnesses}
For one record, activate 1 and prepare a replacement. Intent-only retirement can
then remove 0 while the committed live reference is still 0. A crash clears its
volatile replacement. Alternatively, commit before retiring 0 under live-only:
the live reference is now 1, but the snapshot still requires 0. The guarded rule
rejects that deletion while the snapshot remains in the recovery set.

\subsection*{Different evidence from the two mechanisms}
The model interleaves individual record transitions. The interruption harness
rewraps a complete dataset in one SQL transaction and interrupts an actual
process. The concurrent experiment on the next page adds an admission operation
for an earlier epoch choice; it is not a refinement proof of the closed model.

\newpage
\section{Concurrent admission protocol and schedules}
The coordinator starts two real worker processes, using Python's \texttt{spawn}
method. One retires generation 0 and one attempts to prepare and publish references
to generation 0. Initially the live database and snapshot reference 1, active
generation is 1, both keys exist, and no pending pin exists. Generation 0 is
therefore eligible for retirement at the initial scan.

The writer represents an admitted request carrying an earlier target choice,
as might occur with delayed work or an import. This is an explicit extension
of the finite model: its ordinary prepare operation always selects the active
generation. A current-generation-only implementation can reject old choices;
this experiment instead admits an inactive choice only while its key exists.
It tests reference availability, not whether permitting that choice is a good
profile-migration policy. Already registered pending work is included in both
implementations' scans.

\begin{center}\begin{tabular}{@{}p{2.3cm}p{12.9cm}@{}}
\toprule
Actor & Shared-lock implementation: exact protected interval\\
\midrule
Retirer & Acquire one exclusive lock; inspect committed live references, retained snapshot, registered pending targets, available keys, and active generation; decide eligibility; remove key if permitted; release lock.\\
Writer & Acquire the same lock; validate target-key availability; register a pending pin in SQLite; prepare encrypted replacement envelopes; atomically publish live references and remove the pin in one SQL transaction; release lock.\\
Coordinator & Holds no lifecycle lock; requests actions via IPC and records responses, contention, and inventories. Reopens both copies after the workers finish and checks actual restoration.\\
\bottomrule\end{tabular}\end{center}
Both workers open the same stable lock-file inode. \texttt{flock(LOCK\_EX | LOCK\_NB)}
reports contention without timing sleeps. A blocked action is retried only after
the holder has completed. The unlocked implementation uses the same dependency
scan and key-availability check but omits the shared lock. It can act on a stale
scan. Locking only the deletion would not protect the whole check/delete interval.

Activation and snapshot replacement/release are disabled during these trials.
A deployment enabling them must synchronize their dependency changes too.
The pending registry here is a durable SQLite table; the finite model's pending
work is volatile. No worker failure is injected in the concurrent trials, so
reclaiming abandoned durable pins is not evaluated. The lock is advisory and
local, requiring all participants to cooperate.

\begin{center}\begin{tabular}{@{}lp{9.4cm}@{}}
\toprule
Schedule & Requested order (C: check; P: prepare/register; W: commit; D: delete)\\
\midrule
\texttt{commit\_in\_gap} & C, P, W, D\\
\texttt{prepare\_in\_gap} & C, P, D, W\\
\texttt{writer\_first} & P, C, W, D\\
\texttt{delete\_first} & C, D, P\\
\bottomrule\end{tabular}\end{center}
For a blocked operation, the requested unsafe order is prevented: the holder
finishes first and the blocked operation is retried. Four schedules times two
protocols times six datasets produce 48 cases. Each runs once. This is controlled
interleaving coverage, not exhaustive thread scheduling or a statistical sample.

\newpage
\section{Observed concurrent results and traces}
\begin{center}\begin{tabular}{@{}lrrrr@{}}
\toprule
Schedule & Unlocked L/B & Locked L/B & Locked rejected & Locked committed\\
\midrule
Commit in gap & 6/0 & 0/0 & 6 & 0\\
Prepare in gap & 6/0 & 0/0 & 6 & 0\\
Writer first & 0/0 & 0/0 & 0 & 6\\
Delete first & 0/0 & 0/0 & 6 & 0\\
\bottomrule\end{tabular}\end{center}
L/B counts failed live/backup restores among six datasets per cell. All 24
protected cases restore both copies. In the unlocked implementation, the two
gap schedules fail live restoration in all twelve cases. Both snapshots always
restore because they keep generation 1. When the writer's registered pin exists
before the scan, even the unlocked guard sees it and retains the key.

The following is the observed \texttt{prepare\_in\_gap} trace for one 512-byte
record. Vectors are abbreviated by their single generation; $K$ is the available
key set and $P=-$ means no pending entry. These tables summarize the actual
\texttt{concurrency\_traces.json} events.
\begin{center}\begin{tabular}{@{}rlllll@{}}
\toprule
Step & Unlocked observation & Live & Backup & Pending & Keys\\
\midrule
0 & Initial & 1 & 1 & -- & 0,1\\
1 & Retirer checks: eligible & 1 & 1 & -- & 0,1\\
2 & Writer registers and prepares & 1 & 1 & 0 & 0,1\\
3 & Retirer deletes from stale decision & 1 & 1 & 0 & 1\\
4 & Writer commits and clears pin & 0 & 1 & -- & 1\\
5 & Restore: live fails; backup succeeds & 0 & 1 & -- & 1\\
\bottomrule\end{tabular}\end{center}
\begin{center}\begin{tabular}{@{}rlllll@{}}
\toprule
Step & Shared-lock observation & Live & Backup & Pending & Keys\\
\midrule
0 & Initial & 1 & 1 & -- & 0,1\\
1 & Retirer locks and checks: eligible & 1 & 1 & -- & 0,1\\
2 & Writer's lock attempt is blocked & 1 & 1 & -- & 0,1\\
3 & Retirer deletes and unlocks & 1 & 1 & -- & 1\\
4 & Writer retries: target key absent & 1 & 1 & -- & 1\\
5 & Restore: both copies succeed & 1 & 1 & -- & 1\\
\bottomrule\end{tabular}\end{center}
The raw traces additionally include stop acknowledgments, every inspected
inventory, lock ownership, cached eligibility, pre-deletion inventory, and exact
per-case restore counts. Sequence numbers are coordinator event order, not
latency measurements. Worker command timeouts and nonzero exits abort a run.

\subsection*{Interpretation and limits}
When retirement wins the shared lock, the delayed writer is safely rejected.
When the writer wins, it commits while holding the lock; retirement later sees
the committed dependency and declines deletion. This preserves availability
for these schedules. It does not guarantee eventual retirement or target-profile
completion. Uncooperative writers, distributed key stores, changing snapshot
inventories, worker crashes while holding pins, and all other schedules remain
outside the concurrent evaluation.

\newpage
\section{File-to-claim map and validation}
Paths below are relative to the unzipped artifact root.
\begin{center}\begin{tabular}{@{}p{6.6cm}p{8.7cm}@{}}
\toprule
File & Paper table/claim supported\\
\midrule
\texttt{results/model\_summary.csv} & Table II, including largest-bound state counts.\\
\texttt{results/model\_results.json} & All eighteen configurations and guarded totals.\\
\texttt{results/model\_traces.json} & Shortest missing-key and profile-gap witnesses.\\
\texttt{results/crash\_trials.csv} & Table III: 126 individual interruption/completion cases.\\
\texttt{results/crash\_summary.json} & Table III's policy aggregates.\\
\texttt{results/concurrency\_trials.csv} & Table IV: 48 observed concurrent cases.\\
\texttt{results/concurrency\_summary.json} & Table IV's protocol/schedule aggregates.\\
\texttt{results/concurrency\_traces.json} & Ordered observations and restores for every concurrent case.\\
\texttt{results/*\_manifest.json} & Runtime versions, source hashes, configurations, and scope.\\
\texttt{results/execution\_log.txt} & Successful regression and experiment console output.\\
\texttt{experiment/verify\_results.py} & Checks recorded source hashes, totals, and agreement among result files.\\
\texttt{VALIDATION.md} & Claim corrections, reference audit, and review-response mapping.\\
\texttt{paper/paper.tex}, \texttt{paper.pdf} & Current manuscript with repository URL.\\
\texttt{paper/supplement.tex}, \texttt{supplement.pdf} & This guide and its source (both under \texttt{paper/}).\\
\bottomrule\end{tabular}\end{center}
Tables I and V are source comparison and illustrative acceptance obligations,
not measured outcomes. The figure explains the interruption sequence; it is not
a performance plot. All tables are self-contained in the main manuscript.

\subsection*{Expected reproduction totals and failure conditions}
Nine tests pass: the six original model checks plus late-dependency races,
registered-pending guard coverage, and the four protected schedules. Exploration
produces eighteen configurations. The interruption harness produces 126 cases
(108 terminations), with policy live/backup failures 12/30, 0/12, and 0/0.
The concurrent harness produces 48 cases, with unlocked live failures 12/24,
protected live failures 0/24, and no backup failure. These totals describe
executed reference runs and are checked only after observations are generated.

A timeout, failed worker, unexpected exit, authentication error, wrong row count,
plaintext mismatch, dangling pin, or inventory/restore disagreement is an error.
A protected missing-key state or disagreement between raw files and summaries
also fails validation. Negative-control failures are intentional evidence and
must not be silently converted into successful restores.

\subsection*{Integrity, build, and provenance}
\begin{Verbatim}[fontsize=\small]
sha256sum --check SHA256SUMS
cd paper
latexmk -pdf -interaction=nonstopmode -halt-on-error paper.tex
latexmk -pdf -interaction=nonstopmode -halt-on-error supplement.tex
\end{Verbatim}
The unchanged IEEEtran class is included with its license notices; a TeX Live
installation with TikZ and Latin Modern builds the documents. Rerunning can
change manifests (including the completion timestamp), and rebuilding PDFs can
change embedded dates. Preserve the delivered package before such operations.
OpenAI ChatGPT assisted with the code, model, diagram, text, and validation.
Results were obtained by executing the supplied programs. Author review and
adoption remain necessary; this artifact does not certify a production system.
\end{document}
